\subsection{Deleting a singleton adjacent to a contractible corner}

\begin{lemma}[Proved: the contractible incoming comma]
\label{CP2:A-linear-contractible-singleton}
Let $C$ have full linear support on $0<1<2$, and suppose
$C(0,1)$ has one element.  If $BC(0,0)$ is contractible,
the full-subcategory inclusion
\[
 C|_{\{0,2\}}\longrightarrow C
\]
is a homotopy equivalence on classifying spaces.
There is a dual statement when $C(1,2)$ is singleton and
$BC(2,2)$ is contractible.
\end{lemma}
\begin{proof}
Put $D=C|_{\{0,2\}}$ and write $j:D\hookrightarrow C$.
For a retained object $z$, the comma category $(j\downarrow z)$
has terminal object $(z,1_z)$.  The only object of
$(j\downarrow1)$ is the unique arrow $u:0\to1$: there is
no arrow from $2$ to $1$.  Every $a\in C(0,0)$ gives one
endomorphism of this comma object, since $ua=u$; these are all
its morphisms, with their original multiplication.  Thus this
comma category is literally the one-object category $C(0,0)$.
All these comma categories have contractible classifying spaces.
Quillen's Theorem~A proves the claim.  Taking opposite categories
gives the dual assertion.
\end{proof}

\begin{corollary}[Proved: a small initial singleton at eighteen arrows]
\label{CP2:A-eighteen-small-initial-singleton}
A full-linear three-object category with eighteen arrows,
corner orders $(a,b,c)$ and cross-hom sizes $(x,y,z)$, cannot
have classifying space homotopy equivalent to $\mathbb{CP}^2$
if
\[
 a\leq4,\qquad x=1,\qquad c\leq11.
\]
The opposite condition $c\leq4$, $y=1$, $a\leq11$ also
excludes the category.
\end{corollary}
\begin{proof}
Let $M=C(0,0)$ and let $K$ be its minimal ideal.  By
Lemma~\ref{CP2:A-small-minimal-ideal}, $BM$ is homotopy
equivalent to $BG$ for a finite group $G$ when $|M|\leq4$.
If $G$ is trivial, the preceding lemma deletes object $1$
by an actual homotopy equivalence.  This removes its corner
and both adjacent hom-sets, hence at least three arrows.
It would give a model with at most fifteen arrows,
contradicting Theorem~\ref{CP2:A-global-lower-eighteen}.

Suppose $G$ is nontrivial.  If $M$ itself is a group there
is nothing to prove about its projectivity.  Otherwise
$1_M\notin K$, and the Rees cardinality is
$|K|=rc|G|$, where $r,c$ are its row and column counts.
If $rc\geq2$, then $|M|\geq2\cdot2+1=5$, a contradiction.
Consequently $K$ itself is a group.  Write $e$ for its identity
and put $\varepsilon=|G|^{-1}\sum_{g\in G}g\in\mathbb QM$.
For every $m\in M$, left multiplication on $G$ is the
permutation given by $me\in G$.  Hence $m\varepsilon=\varepsilon$.
Also $\varepsilon^2=\varepsilon$, so the left ideal
$\mathbb QM\varepsilon$ is precisely the trivial
one-dimensional left module and is projective.  The same
averaging argument applies when $M$ is a group.

The singleton adjacent act has one orbit.  Corollary
\ref{CP2:A-eleven-projective-singleton} therefore annihilates
all rational degree-two products, since the final corner
has order at most eleven and the assumed projective-plane
homotopy type supplies bounded rational homology.  This
contradicts the nonzero square on $\mathbb{CP}^2$.
Opposition proves the last assertion.
\end{proof}
